% =============================================================================
% Chapter 6: Results
% Findings only. All interpretation is in Chapter 7.
% Every figure block is copied verbatim from figure_captions.tex and every table
% from tables.tex or dashboard_tables.tex, so the numbers cannot drift.
% Results are reported as actual run counts with their eligible population;
% rates per 100,000 appear only inside an SPI definition, never as a result.
% =============================================================================

\chapter{Results}
\label{ch:results}

This chapter reports what the register of Chapter~\ref{ch:methodology} does on
the 94,816 analysed runs of Chapter~\ref{ch:experiment}. It states the findings
and does not interpret them; the interpretation follows in
Chapter~\ref{ch:discussion}. Every result is given as a number of runs together
with the eligible population of the chain concerned.

Section~\ref{sec:results_collisions} gives the collisions each chain produced.
Section~\ref{sec:results_violations} sets the violations of the leading SPIs
against those collisions, and Section~\ref{sec:results_leadtime} reports which
collisions were preceded by a violation and by how long.
Section~\ref{sec:results_statistics} gives the detection statistics.
Sections~\ref{sec:results_gate} and~\ref{sec:results_poolability} report two
properties of the evaluation itself, namely the effect of the oncoming conflict
gate and the transferability of the metric thresholds between campaigns.
Section~\ref{sec:results_dashboard} gives the state of the register as the
monitoring dashboard presents it.

% =============================================================================
\section{Collisions per Conflict Chain}
\label{sec:results_collisions}

The collisions are the reference events of the whole evaluation, so they are
stated first. Table~\ref{tab:collisions_per_chain} gives them per chain and per
campaign under the collision definition of
Section~\ref{sec:exp_collision}.

% ---------------------------------------------------------------------------
\begin{table}[htbp]
    \centering
    \caption{Collisions per conflict chain. The table counts the collision runs
    of each chain per campaign under the collision definition of
    Section~\ref{sec:exp_collision}.}
    \label{tab:collisions_per_chain}
    \small
    \begin{tabular}{@{}lrrrr@{}}
        \toprule
        {\bfseries Chain} & {\bfseries Campaign 1} & {\bfseries Campaign 2} &
        {\bfseries Campaign 3} & {\bfseries Total} \\
        \midrule
        Oncoming vehicle (ONC)  & 540 & 0 & 0 & 540 \\
        Pedestrian (PED)        & 814 & 1 & 0 & 815 \\
        Following vehicle (REAR)& 101 & 0 & 0 & 101 \\
        \bottomrule
    \end{tabular}
\end{table}

All 540 oncoming-vehicle collisions and all 101 following-vehicle collisions
occur in Campaign~1. Of the 815 pedestrian collisions, 814 occur in Campaign~1
and one in Campaign~2. The 75,509 runs of Campaigns~2 and~3 together therefore
contain a single collision across all three chains.

% =============================================================================
\section{Violations and Collisions}
\label{sec:results_violations}

The question this section answers is how often each leading SPI was violated,
set against how often its chain produced a collision, and whether that relation
survives the change from the provoked configuration to realistic driving.
Figure~\ref{fig:6_1_violations_collisions} and
Table~\ref{tab:violations_collisions} give the counts.

% ---------------------------------------------------------------------------
% Copied verbatim from figure_captions.tex
\begin{figure}[htbp]
\centering
\includegraphics[width=16cm]{fig_6_1_violations_collisions.pdf}
\caption{Leading SPI violations and collisions per conflict chain. Bars give the number of eligible runs in which each leading SPI was violated and the number of collision runs of the chain (hatched), split into Campaign 1 and Campaigns 2 and 3; the metric thresholds are calibrated on Campaign 1 and applied unchanged to all 94,816 runs.}
\label{fig:6_1_violations_collisions}
\end{figure}

% ---------------------------------------------------------------------------
% Copied verbatim from tables.tex
\begin{table}[htbp]
\centering
\caption{Violations and collisions. The table gives the number of runs per configuration in which each leading SPI was violated, and the collision runs of its chain.}
\label{tab:violations_collisions}
\small
\begin{tabular}{@{}llrrrrr@{}}
\toprule
{\bfseries Indicator} & {\bfseries Chain} & {\bfseries \begin{tabular}[b]{@{}c@{}}Eligible\\runs\end{tabular}} & {\bfseries \begin{tabular}[b]{@{}c@{}}Violations\\Campaign 1\end{tabular}} & {\bfseries \begin{tabular}[b]{@{}c@{}}Violations\\Campaigns 2, 3\end{tabular}} & {\bfseries \begin{tabular}[b]{@{}c@{}}Violations\\total\end{tabular}} & {\bfseries \begin{tabular}[b]{@{}c@{}}Collisions\\C1 / C2, 3\end{tabular}} \\
\midrule
GAP\_TTC & ONC & 80,455 & 2,070 & 5,624 & 7,694 & 540 / 0 \\
REQ\_DECEL & ONC & 80,455 & 1,031 & 1,218 & 2,249 & 540 / 0 \\
AEB\_ONC & ONC & 80,455 & 629 & 753 & 1,382 & 540 / 0 \\
PED\_CPA & PED & 37,309 & 1,268 & 527 & 1,795 & 814 / 1 \\
AEB\_PED & PED & 37,309 & 851 & 657 & 1,508 & 814 / 1 \\
REAR\_TTC & REAR & 89,882 & 2,243 & 9,841 & 12,084 & 101 / 0 \\
\bottomrule
\end{tabular}
\end{table}

The eligible populations differ between the two configurations. For the oncoming
chain, 15,362 of the 80,455 eligible runs belong to Campaign~1 and 65,093 to
Campaigns~2 and~3; for the pedestrian chain, 8,424 and 28,885; for the
following-vehicle chain, 14,373 and 75,509.

In Campaign~1 the violation counts and the collision counts are of the same
order. The oncoming chain records 2,070 GAP\_TTC violations against 540
collisions, which is 3.8 violations per collision, 1,031 REQ\_DECEL violations at
1.9 per collision and 629 AEB\_ONC events at 1.2 per collision. The pedestrian
chain records 1,268 PED\_CPA violations against 814 collisions, or 1.6 per
collision, and 851 AEB\_PED events at 1.0 per collision. The following-vehicle
chain records 2,243 REAR\_TTC violations against 101 collisions, or 22.2 per
collision.

In Campaigns~2 and~3 the relation changes. The 75,509 runs contain one collision,
while every leading SPI continues to be violated: 5,624 times for GAP\_TTC, 1,218
for REQ\_DECEL, 753 for AEB\_ONC, 527 for PED\_CPA, 657 for AEB\_PED and 9,841
for REAR\_TTC. For the pedestrian chain this gives 527 PED\_CPA violations and
657 AEB\_PED events per collision. For the chains without any collision no ratio
is defined; using the 95\,\% Wilson upper bound of a zero collision count gives
at least 1,464 violations per collision for GAP\_TTC, at least 317 for
REQ\_DECEL, at least 196 for AEB\_ONC and at least 2,562 for REAR\_TTC.

% =============================================================================
\section{Detection and Lead Time}
\label{sec:results_leadtime}

A violation is only useful if it precedes the collision it anticipates, and the
time by which it precedes it bounds any response. This section reports both.
Figure~\ref{fig:6_2_detection} gives the share of collision runs with an earlier
violation, Figure~\ref{fig:6_3_lead_time} the distribution of the lead time, and
Table~\ref{tab:lead_time} both as counts. Lead time is measured from the first
violation in the run to the collision.

% ---------------------------------------------------------------------------
% Copied verbatim from figure_captions.tex
\begin{figure}[htbp]
\centering
\includegraphics[width=16cm]{fig_6_2_detection.pdf}
\caption{Collision runs with a violation before the collision. For each leading SPI, the bar divides the collision runs of its chain into runs in which the SPI was violated before the collision and runs without such a violation, over all 94,816 analysed runs.}
\label{fig:6_2_detection}
\end{figure}

% ---------------------------------------------------------------------------
% Copied verbatim from figure_captions.tex
\begin{figure}[htbp]
\centering
\includegraphics[width=16cm]{fig_6_3_lead_time.pdf}
\caption{Lead time of the leading SPIs. Markers give the median and bars the interquartile range of the time from the first violation to the collision in collision runs with an earlier violation; the numbers at the right count collisions warned at least 1\,s ahead. AEB\_PED is omitted because its lead time does not measure a warning.}
\label{fig:6_3_lead_time}
\end{figure}

% ---------------------------------------------------------------------------
% Copied verbatim from tables.tex
\begin{table}[htbp]
\centering
\caption{Detection and lead time. The table counts collision runs with an earlier violation and gives the time from the first violation to the collision.}
\label{tab:lead_time}
\small
\begin{tabular}{@{}llrrrr@{}}
\toprule
{\bfseries Indicator} & {\bfseries Chain} & {\bfseries \begin{tabular}[b]{@{}c@{}}Violation before\\collision\end{tabular}} & {\bfseries \begin{tabular}[b]{@{}c@{}}Median\\lead\end{tabular}} & {\bfseries Q1 to Q3} & {\bfseries \begin{tabular}[b]{@{}c@{}}Warned at least\\1\,s ahead\end{tabular}} \\
\midrule
GAP\_TTC & ONC & 527 of 540 & 0.60\,s & 0.55 to 0.70\,s & 38 of 540 \\
REQ\_DECEL & ONC & 317 of 540 & 0.10\,s & 0.10 to 0.20\,s & 1 of 540 \\
AEB\_ONC & ONC & 286 of 540 & 0.75\,s & 0.20 to 0.95\,s & 57 of 540 \\
PED\_CPA & PED & 790 of 815 & 0.50\,s & 0.40 to 0.55\,s & 2 of 815 \\
AEB\_PED & PED & 110 of 815 & 5.15\,s & 4.95 to 6.24\,s & not a warning time \\
REAR\_TTC & REAR & 101 of 101 & 1.15\,s & 1.05 to 3.30\,s & 87 of 101 \\
\bottomrule
\end{tabular}
\par\smallskip\footnotesize\raggedright Lead time is measured from the first violation in the run.
\end{table}

The two kinematic indicators that watch the first branch of their chain precede
nearly every collision of that chain: GAP\_TTC 527 of 540 oncoming-vehicle
collisions, PED\_CPA 790 of 815 pedestrian collisions and REAR\_TTC all 101
following-vehicle collisions. REQ\_DECEL precedes 317 of 540 and AEB\_ONC 286 of
540 oncoming-vehicle collisions. AEB\_PED precedes 110 of 815 pedestrian
collisions.

The lead times separate the following conflict from the two crossing conflicts.
REAR\_TTC has a median lead of 1.15\,s and warns at least one second ahead in 87
of the 101 following-vehicle collisions. The three oncoming-vehicle indicators
have median leads of 0.60\,s, 0.10\,s and 0.75\,s and warn at least one second
ahead in 38, one and 57 of the 540 collisions respectively. PED\_CPA has a median
lead of 0.50\,s and warns at least one second ahead in two of the 815 pedestrian
collisions.

The AEB\_PED row is not comparable with the others. Only 110 of the 815
pedestrian collisions follow an emergency braking event at all, and in those 110
runs the event lies between 4.95\,s and 6.24\,s before the collision, with a
median of 5.15\,s. Because lead time is measured from the first violation
anywhere in the run, this interval is not established to be a warning time for
the collision that follows it, and it is not reported as one. The point is taken
up in Section~\ref{sec:disc_aebped}.

% =============================================================================
\section{Detection Statistics}
\label{sec:results_statistics}

The counts of the previous section describe each indicator from one side at a
time. This section combines them into the classification of
Section~\ref{sec:exp_criteria} and reports precision, recall, false-alarm rate
and relative risk with their intervals. Table~\ref{tab:confusion_counts} gives
the classification of every eligible run and
Table~\ref{tab:detection_statistics} the measures derived from it.

% ---------------------------------------------------------------------------
% Copied verbatim from tables.tex
\begin{table}[htbp]
\centering
\caption{Confusion counts. The table classifies every eligible run by whether the SPI was violated before any collision and whether a collision of the chain occurred.}
\label{tab:confusion_counts}
\small
\begin{tabular}{@{}llrrrr@{}}
\toprule
{\bfseries Indicator} & {\bfseries Chain} & {\bfseries TP} & {\bfseries FP} & {\bfseries FN} & {\bfseries TN} \\
\midrule
GAP\_TTC & ONC & 527 & 7,167 & 13 & 72,748 \\
REQ\_DECEL & ONC & 317 & 1,932 & 223 & 77,983 \\
AEB\_ONC & ONC & 286 & 1,096 & 254 & 78,819 \\
PED\_CPA & PED & 790 & 1,005 & 25 & 35,489 \\
AEB\_PED & PED & 110 & 1,398 & 705 & 35,096 \\
REAR\_TTC & REAR & 101 & 11,983 & 0 & 77,798 \\
\bottomrule
\end{tabular}
\par\smallskip\footnotesize\raggedright TP: violation before a collision of the chain; FP: violation, no collision; FN: collision without an earlier violation; TN: neither.
\end{table}

% ---------------------------------------------------------------------------
% Copied verbatim from tables.tex
\begin{table}[htbp]
\centering
\caption{Detection statistics. Precision, recall, false-alarm rate and relative risk of each leading SPI over all eligible runs, with 95\,\% intervals.}
\label{tab:detection_statistics}
\footnotesize
\setlength{\tabcolsep}{4pt}
\begin{tabular}{@{}lrrrr@{}}
\toprule
{\bfseries Indicator} & {\bfseries \begin{tabular}[b]{@{}c@{}}Precision\\in \%\end{tabular}} & {\bfseries \begin{tabular}[b]{@{}c@{}}Recall\\in \%\end{tabular}} & {\bfseries \begin{tabular}[b]{@{}c@{}}False-alarm\\rate in \%\end{tabular}} & {\bfseries Relative risk} \\
\midrule
GAP\_TTC & 6.8 [6.3, 7.4] & 97.6 [95.9, 98.6] & 8.97 & 383.4 [221.2, 664.3] \\
REQ\_DECEL & 14.1 [12.7, 15.6] & 58.7 [54.5, 62.8] & 2.42 & 49.4 [41.9, 58.4] \\
AEB\_ONC & 20.7 [18.6, 22.9] & 53.0 [48.7, 57.1] & 1.37 & 64.4 [54.9, 75.6] \\
PED\_CPA & 44.0 [41.7, 46.3] & 96.9 [95.5, 97.9] & 2.75 & 625.2 [421.1, 928.3] \\
AEB\_PED & 7.3 [6.1, 8.7] & 13.5 [11.3, 16.0] & 3.83 & 3.7 [3.1, 4.5] \\
REAR\_TTC & 0.8 [0.7, 1.0] & 100.0 [96.3, 100.0] & 13.35 & $\geq$\,1,306.8 [81.2, 21,035.8] \\
\bottomrule
\end{tabular}
\par\smallskip\footnotesize\raggedright Precision and recall with Wilson score intervals; relative risk with the Katz log interval. REAR\_TTC: every collision was preceded by a violation, so the relative risk is a Haldane-corrected lower bound.
\end{table}

Every leading SPI has a relative risk whose 95\,\% interval excludes unity. The
point estimates are 383.4 for GAP\_TTC, 49.4 for REQ\_DECEL, 64.4 for AEB\_ONC,
625.2 for PED\_CPA and 3.7 for AEB\_PED; REAR\_TTC records no collision without
an earlier violation, so its value of 1,306.8 is a lower bound and its interval,
from 81.2 to 21,035.8, is correspondingly wide.

The pooled precision values rest on all three campaigns together, and
Campaigns~2 and~3 contribute violations while contributing one collision. The
pooled precision is therefore lower than the precision within Campaign~1 for
every indicator, and the two values are reported together.
Table~\ref{tab:statistics_per_campaign}, placed in Appendix~\ref{app:statistics},
gives the full per-campaign breakdown. Within Campaign~1 the precision is
25.5\,\% for GAP\_TTC against 6.8\,\% pooled, 30.7\,\% against 14.1\,\% for
REQ\_DECEL, 45.5\,\% against 20.7\,\% for AEB\_ONC, 62.2\,\% against 44.0\,\%
for PED\_CPA, 12.8\,\% against 7.3\,\% for AEB\_PED and 4.5\,\% against 0.8\,\%
for REAR\_TTC. The relative risks within Campaign~1 are 260.3 for GAP\_TTC, 19.8
for REQ\_DECEL, 26.4 for AEB\_ONC, 178.1 for PED\_CPA, 1.4 for AEB\_PED and at
least 1,097.4 for REAR\_TTC.

% =============================================================================
\section{Effect of the Oncoming Conflict Gate}
\label{sec:results_gate}

Step~5 of the derivation requires that an indicator be computed only where its
conflict can arise. For the oncoming chain this is implemented by the conflict
gate defined in Chapter~\ref{ch:methodology}, and its effect is isolated by scoring the
three oncoming-vehicle indicators twice on the same runs, once with and once
without the gate. Table~\ref{tab:conflict_gate} gives the comparison.

% ---------------------------------------------------------------------------
% Copied verbatim from tables.tex
\begin{table}[htbp]
\centering
\caption{Oncoming conflict gate. The table compares the oncoming chain without and with the conflict gate.}
\label{tab:conflict_gate}
\small
\begin{tabular}{@{}lrrr@{}}
\toprule
{\bfseries Indicator} & {\bfseries Eligible runs} & {\bfseries Violations} & {\bfseries ONC collisions} \\
\midrule
GAP\_TTC & 94,816 / 80,455 & 10,360 / 7,694 & 540 / 540 \\
REQ\_DECEL & 94,816 / 80,455 & 3,146 / 2,249 & 540 / 540 \\
AEB\_ONC & 94,816 / 80,455 & 1,490 / 1,382 & 540 / 540 \\
\bottomrule
\end{tabular}
\par\smallskip\footnotesize\raggedright Values given as without gate / with gate.
\end{table}

The gate removes 14,361 runs from the eligible population of the oncoming chain
and all 540 oncoming-vehicle collisions remain eligible. The violation counts
fall from 10,360 to 7,694 for GAP\_TTC, from 3,146 to 2,249 for REQ\_DECEL and
from 1,490 to 1,382 for AEB\_ONC.

% =============================================================================
\section{Poolability of the Metric Thresholds}
\label{sec:results_poolability}

The metric thresholds are calibrated on Campaign~1 and applied to all three
campaigns, which is defensible only if the underlying distributions are close
enough for one value to stand for all of them. To test this, each threshold was
recalibrated separately on each campaign and the pairwise relative shift
examined. Table~\ref{tab:poolability} gives the result; pooling is refused above
a relative shift of 15\,\%.

% ---------------------------------------------------------------------------
% Copied verbatim from tables.tex
\begin{table}[htbp]
\centering
\caption{Threshold poolability. Metric thresholds calibrated separately per campaign and their relative shift.}
\label{tab:poolability}
\small
\begin{tabular}{@{}lrrrrrr@{}}
\toprule
{\bfseries Indicator} & {\bfseries $\theta$ C1} & {\bfseries $\theta$ C2} & {\bfseries $\theta$ C3} & {\bfseries Shift 1 vs 2} & {\bfseries Shift 1 vs 3} & {\bfseries Shift 2 vs 3} \\
\midrule
GAP\_TTC & 0.913 & 0.927 & 0.926 & 1.55\,\% & 1.50\,\% & 0.05\,\% \\
REQ\_DECEL & 7.549 & 5.924 & 5.924 & 21.52\,\%, no & 21.52\,\%, no & 0.00\,\% \\
PED\_CPA & 1.174 & 1.619 & 1.617 & 37.90\,\%, no & 37.74\,\%, no & 0.12\,\% \\
REAR\_TTC & 2.679 & 2.517 & 2.500 & 6.08\,\% & 6.71\,\% & 0.67\,\% \\
\bottomrule
\end{tabular}
\par\smallskip\footnotesize\raggedright Pooling is refused above 15\,\% relative shift (marked no).
\end{table}

Campaigns~2 and~3 agree on every metric threshold to within 0.67\,\%, so
Campaign~3 is a statistical replicate of Campaign~2. Between Campaign~1 and the
two realistic-driving campaigns, the thresholds of GAP\_TTC and REAR\_TTC remain
within the criterion, at shifts of 1.55\,\% and 1.50\,\% and of 6.08\,\% and
6.71\,\%. The thresholds of REQ\_DECEL and PED\_CPA do not: REQ\_DECEL shifts by
21.52\,\% against both campaigns and PED\_CPA by 37.90\,\% and 37.74\,\%.

% =============================================================================
\section{The Register under Monitoring}
\label{sec:results_dashboard}

The preceding sections evaluate the register after the fact. A monitoring process
reads it while the runs accumulate, and the form in which it does so is part of
the result. The nine SPIs were exported to a database and presented in a Grafana
dashboard; Figure~\ref{fig:6_5_dashboard} shows it. The dashboard is not deployed
on a vehicle.

% ---------------------------------------------------------------------------
\begin{figure}[htbp]
    \centering
    \includegraphics[width=16cm]{fig_6_5_dashboard.png}
    \caption{The SPI register under monitoring. The dashboard presents each
    indicator with its metric threshold, the baseline violations of
    Campaigns~2 and~3, a detection scorecard per chain and a funnel from
    eligible runs to collisions, over the same results as the preceding
    sections.}
    \label{fig:6_5_dashboard}
\end{figure}

The dashboard presents the register itself, a detection scorecard per chain, a
funnel from eligible runs to collisions, and tables of violations and collisions
per operating condition.
Tables~\ref{tab:dashboard_register} to~\ref{tab:dashboard_funnel} give the three
of these that bear on the register as a whole;
Table~\ref{tab:dashboard_odd}, in Appendix~\ref{app:statistics}, gives the
operating-condition breakdown.

% ---------------------------------------------------------------------------
% Trimmed from dashboard_tables.tex (v4: window columns removed)
\begin{table}[htbp]
\centering
\caption{SPI register as monitored. Identifier, type and metric threshold of each SPI, together with its baseline violations in Campaigns~2 and~3.}
\label{tab:dashboard_register}
\small
\begin{tabular}{@{}llrr@{}}
\toprule
{\bfseries SPI} & {\bfseries Type} & {\bfseries $\theta$} & {\bfseries \begin{tabular}[b]{@{}c@{}}Baseline\\(Campaigns 2, 3)\end{tabular}} \\
\midrule
LTI-ONC-P1-REQ\_DECEL & leading & 7.549\,m/s\textsuperscript{2} & 1,218 of 65,093 \\
LTI-ONC-P2-GAP\_TTC & leading & 0.913\,s & 5,624 of 65,093 \\
LTI-ONC-P3-AEB\_ONC & leading & -- & 753 of 65,093 \\
LTI-PED-P3-AEB\_PED & leading & -- & 657 of 28,885 \\
LTI-PED-S1-PED\_CPA & leading & 1.174\,s & 527 of 28,885 \\
LTI-REAR-P1-REAR\_TTC & leading & 2.679\,s & 9,841 of 75,509 \\
LTI-ONC-O1-COLLISION & lagging & -- & 0 of 65,093 \\
LTI-PED-O1-COLLISION & lagging & -- & 1 of 28,885 \\
LTI-REAR-O1-COLLISION & lagging & -- & 0 of 75,509 \\
\bottomrule
\end{tabular}
\par\smallskip\footnotesize\raggedright Baseline: eligible runs of Campaigns~2 and~3 with a violation of the SPI; lagging rows count collisions. The SPI threshold T of each indicator is derived from this baseline as defined in Equation~(\ref{eq:exp_spithreshold}).
\end{table}

The register holds nine SPIs, six leading and three lagging, each identified by
use case, chain, tier and signal. In the baseline of Campaigns~2 and~3 the six
leading SPIs record between 527 and 9,841 violating runs, whereas the three
lagging SPIs record one collision in total, which belongs to the pedestrian
chain.

% ---------------------------------------------------------------------------
% Copied verbatim from dashboard_tables.tex
\begin{table}[htbp]
\centering
\caption{Chain detection scorecard. Collision runs of each chain with a violation before the collision, for each SPI and for any SPI of the chain.}
\label{tab:dashboard_scorecard}
\small
\begin{tabular}{@{}llrrrr@{}}
\toprule
{\bfseries Chain} & {\bfseries Indicator} & {\bfseries \begin{tabular}[b]{@{}c@{}}Collision\\runs\end{tabular}} & {\bfseries \begin{tabular}[b]{@{}c@{}}Violation before\\collision\end{tabular}} & {\bfseries \begin{tabular}[b]{@{}c@{}}Lead at\\least 1\,s\end{tabular}} & {\bfseries \begin{tabular}[b]{@{}c@{}}Median\\lead\end{tabular}} \\
\midrule
ONC & GAP\_TTC & 540 & 527 & 38 & 0.60\,s \\
ONC & REQ\_DECEL & 540 & 317 & 1 & 0.10\,s \\
ONC & AEB\_ONC & 540 & 286 & 57 & 0.75\,s \\
ONC & any ONC & 540 & 538 & 92 & 0.70\,s \\
PED & PED\_CPA & 815 & 790 & 2 & 0.50\,s \\
PED & AEB\_PED & 815 & 110 & 108 & 5.15\,s \\
PED & any PED & 815 & 791 & 110 & 0.55\,s \\
REAR & REAR\_TTC & 101 & 101 & 87 & 1.15\,s \\
REAR & any REAR & 101 & 101 & 87 & 1.15\,s \\
\bottomrule
\end{tabular}
\par\smallskip\footnotesize\raggedright \emph{any}: earliest violation among the chain's indicators.
\end{table}

Taking the earliest violation among the indicators of a chain raises the
detection beyond that of any single indicator for the oncoming and pedestrian
chains: 538 of 540 oncoming-vehicle collisions and 791 of 815 pedestrian
collisions are preceded by a violation of at least one indicator of the chain,
against 527 and 790 for the best single indicator. The number warned at least one
second ahead rises from 57 to 92 for the oncoming chain and from two to 110 for
the pedestrian chain.

% ---------------------------------------------------------------------------
% Copied verbatim from dashboard_tables.tex
\begin{table}[htbp]
\centering
\caption{Leading-to-lagging funnel. Eligible runs of each chain with at least one kinematic leading violation, an emergency braking event and a collision.}
\label{tab:dashboard_funnel}
\small
\begin{tabular}{@{}lrrrr@{}}
\toprule
{\bfseries Chain} & {\bfseries Eligible runs} & {\bfseries \begin{tabular}[b]{@{}c@{}}Leading\\violation\end{tabular}} & {\bfseries \begin{tabular}[b]{@{}c@{}}Emergency\\braking\end{tabular}} & {\bfseries Collision} \\
\midrule
Oncoming vehicle & 80,455 & 9,156 & 1,382 & 540 \\
Pedestrian & 37,309 & 1,795 & 1,508 & 815 \\
Following vehicle & 89,882 & 12,084 & 0 & 101 \\
\bottomrule
\end{tabular}
\end{table}

The funnel narrows at every stage for the oncoming chain, from 80,455 eligible
runs through 9,156 with a kinematic leading violation and 1,382 with an emergency
braking event to 540 collisions. For the pedestrian chain the two leading stages
are of similar size, 1,795 and 1,508, against 815 collisions. The
following-vehicle chain has no reaction tier, so its funnel runs from 89,882
eligible runs through 12,084 leading violations to 101 collisions.